% Compatibility entry point; active chapter source: chapter1_introduction.tex.
\section{Introduction and Project Objectives}
\label{ch:introduction}

Image captioning is the task of describing an image in a short natural-language sentence. A useful caption names important objects and may describe their attributes, actions, or relationships. A caption should also avoid details that are not supported by the image. Current vision-language systems can produce fluent text, but fluent wording alone does not guarantee that the details are correct.

Retrieval-Augmented Generation (RAG) is a design in which a system retrieves information from a separate collection and uses it when generating an answer \citep{lewis2020rag}. For image captioning, the retrieved material could include visually similar training images, captions written for those images, or structured visual metadata. The retrieved examples may guide a caption generator, but they may also introduce incorrect objects or relationships. Their effect must be measured rather than assumed.

\subsection{Objective}
The project asks how retrieval techniques affect caption quality, relevance, and grounding. This report covers Phase 1: the data, metadata, embedding, and retrieval foundation. It does not generate captions or compare RAG methods.

\endinput
\chapter{Introduction and Problem Statement}\label{ch:intro}
Image captioning maps a visual scene to a concise natural-language description. A useful caption must identify salient entities, their attributes, and relevant relations while avoiding details unsupported by the image. Modern vision-language encoders provide transferable representations, but caption generation remains sensitive to data coverage, retrieval quality, decoder behavior, and evaluation design. Retrieval-augmented generation (RAG) introduces external examples or knowledge at inference time, potentially giving a generator access to information beyond its parametric weights \cite{lewis2020rag}.

The project title proposes a comparative study of RAG techniques for generating image captions. The implemented repository currently supplies important groundwork: fixed data splits, five references per image, visual and linguistic analyses, CLIP representations, and train-only similarity indexes. Its present status must be understood as an infrastructure and data-preparation phase, not the completed comparative study. This distinction governs every result reported here.

The report has three aims. First, it describes the code and artifacts that exist, rather than treating the project title as proof of completed experiments. Second, it summarizes verified properties of the dataset and retrieval pipeline. Third, it specifies how the pending generation study can be conducted without implying that proposed metrics or outcomes have already been measured.

\section{Problem Statement}\label{ch:problem}
The intended research question is whether retrieval-augmented captioning methods can improve the quality and factual grounding of image captions relative to a non-retrieval baseline. Answering it requires (i) selecting comparable caption generators and RAG strategies, (ii) retrieving relevant evidence without test leakage, (iii) generating captions under controlled decoding settings, and (iv) evaluating linguistic quality and visual faithfulness on held-out images. The repository has completed only part of the second requirement: it creates train-only image and caption vector indexes and checks image-neighbor retrieval.

The current operational problem is therefore narrower: can the project assemble a stable, audited COCO subset and produce correctly mapped visual/text representations and retrieval structures suitable for later experiments? The recorded evidence supports a positive answer for the implemented Phase 1 checks. It does not establish a causal effect of retrieval on generated captions.

\section{Research questions and scope}
The completed phase addresses data integrity and pipeline feasibility. It does not answer which RAG method performs best, whether RAG reduces hallucinations, or whether improvements are statistically significant. Those are Phase 2 questions. The available retrieval smoke test reports embedding similarity, not reference-caption agreement or factual correctness.

\section{Test-size requirement}
The original assignment's detailed dataset and preprocessing requirements specify 1,000 test images. One occurrence contains the typographical value \code{1,00}; the detailed requirements and data-audit instructions establish 1,000 as intended. The frozen split contains 1,000 test images and is unchanged.
